\chapter{Compactification on a Circle: Momentum and Winding}\label{ch:circle}

Suppose one of the 25 spatial directions of the bosonic string is not a line but a circle of
radius $R$. What does an observer in the remaining 25 dimensions see? For a point particle the
answer has been known since Kaluza and Klein: a tower of particles with masses $|n|/R$,
$n\in\Z$, charged under a gauge field that descends from the higher-dimensional metric. This
chapter derives that answer and then asks the same question of a string. The string gives a
second tower, with masses $|w|R/\ap$, $w\in\Z$, coming from strings that wrap the circle. No
point particle can do that. The result is the mass formula
\begin{equation}\label{eq:circle:preview}
 M^2=\frac{n^2}{R^2}+\frac{w^2R^2}{\ap^2}+\frac{2}{\ap}\big(N+\tilde N-2\big),\qquad
 N-\tilde N=nw ,
\end{equation}
which we derive step by step from the mode expansion of \cref{ch:classical} and the constraints
of \cref{ch:quantum}. A student should care because the rest of Part~II and all of Part~III
start from this formula. T-duality (\cref{ch:tduality}) is the observation that
\eqref{eq:circle:preview} does not change under $R\to\ap/R$, $n\leftrightarrow w$. The Narain
lattice (\cref{ch:narain}) is the set of pairs $(p_L,p_R)$ introduced below. The doubled
coordinates of double field theory (\cref{ch:doubled}) are the variables conjugate to $n$ and
$w$. The ``dual-scale'' principle of \cref{ch:dualscale} is a statement about the first two
terms of \eqref{eq:circle:preview}. The chapter closes with the Lean declarations that encode
the momentum part of \eqref{eq:circle:preview}, and with one theorem, \lean{massForm\_circle},
which proves that two independently written libraries compute the same number.

Conventions are those of the book: $\hbar=c=1$, mostly-plus signature, $\ap$ of dimension
length$^2$, string length $\ell_s=\sqrt{\ap}$, tension $T=1/(2\pi\ap)$, closed strings with
$\sigma\sim\sigma+2\pi$ and $\sigma^\pm=\tau\pm\sigma$. The compact coordinate is $X^{25}$,
written $X$ when no confusion can arise; the non-compact coordinates are $X^\mu$ with
$\mu=0,\dots,24$. The integer $n$ is the momentum number and $w$ the winding number. Tong
writes $m$ for our $w$; Giveon, Porrati and Rabinovici write $m$ as well and use dimensionless
momenta, as explained in \cref{rem:circle:gpr}.

%=====================================================================================
\section{Kaluza--Klein reduction of a field}\label{sec:circle:kk}

\subsection{A scalar field on \texorpdfstring{$\R^{1,D-2}\times S^1$}{R x S1}}

Before the string, take the simplest object that can live on a spacetime with one circular
direction: a free massless real scalar field $\Phi(x^\mu,y)$ in $D$ dimensions, where
$x^\mu$ are coordinates on $\R^{1,D-2}$ and $y\sim y+2\pi R$ is the coordinate on the circle.
Its action is
\begin{equation}\label{eq:circle:scalaraction}
 S=-\frac12\int\dd^{D-1}x\int_0^{2\pi R}\dd y\,
   \Big(\partial_\mu\Phi\,\partial^\mu\Phi+(\partial_y\Phi)^2\Big).
\end{equation}
The field must be single valued on the circle, $\Phi(x,y+2\pi R)=\Phi(x,y)$, so it has a
Fourier series\index{Kaluza--Klein reduction}
\begin{equation}\label{eq:circle:fourier}
 \Phi(x,y)=\sum_{n\in\Z}\Phi_n(x)\,\ee^{\ii ny/R},\qquad \Phi_{-n}=\Phi_n^{*},
\end{equation}
where the second condition makes $\Phi$ real \source{Tong §8.1.1, ll.~11283--11310}. Insert
\eqref{eq:circle:fourier} into \eqref{eq:circle:scalaraction}. The $y$-integral of a product of
two exponentials is
\[
 \int_0^{2\pi R}\dd y\;\ee^{\ii(n+m)y/R}=2\pi R\,\delta_{n+m,0},
\]
so only pairs $(n,-n)$ survive. For the $y$-derivative term,
$\partial_y\Phi=\sum_n(\ii n/R)\Phi_n\ee^{\ii ny/R}$, and the pair $(n,m=-n)$ contributes
$(\ii n/R)(-\ii n/R)\Phi_n\Phi_{-n}=(n^2/R^2)|\Phi_n|^2$. The result is
\begin{equation}\label{eq:circle:kkaction}
 S=-\frac{2\pi R}{2}\sum_{n\in\Z}\int\dd^{D-1}x\,
   \Big(\partial_\mu\Phi_n\,\partial^\mu\Phi_{-n}+\frac{n^2}{R^2}\,|\Phi_n|^2\Big).
\end{equation}
This is the action of infinitely many fields in $D-1$ dimensions. The field $\Phi_0$ is real
and massless. For each $n\ge1$ the pair $\Phi_{\pm n}$ is one complex field whose mass can be
read off from the term without derivatives:
\begin{equation}\label{eq:circle:kkmass}
 M_n^2=\frac{n^2}{R^2}
\end{equation}
\TL\ \source{Tong §8.1.1, eq.~(8.3), ll.~11311--11335}\index{Kaluza--Klein tower}.

There is a one-line way to remember \eqref{eq:circle:kkmass}. A massless particle in $D$
dimensions has $p_Mp^M=0$. Split the momentum as $p_M=(p_\mu,p_y)$. An observer who cannot
resolve the circle measures the $(D-1)$-dimensional mass $M^2=-p_\mu p^\mu=p_y^2$. The
wavefunction $\ee^{\ii p_yy}$ is single valued only if $p_y=n/R$. Momentum along a hidden
direction looks like mass.

\begin{physicsbox}[Physical picture: why a small circle is invisible]
The modes with $n\neq0$ cost an energy of at least $1/R$. An experiment with energy
$E\ll1/R$ cannot produce them, and sees only $\Phi_0(x)$: a field that does not depend on
$y$ at all. In this sense the extra dimension disappears at low energy; equivalently, a
probe of wavelength $\lambda\gg R$ cannot resolve the circle. In ordinary units the first
Kaluza--Klein mass is $M_1c^2=\hbar c/R$ with $\hbar c\simeq197.3\ \mathrm{MeV\,fm}$. A circle
of radius $R=10^{-19}\,\mathrm{m}=10^{-4}\,\mathrm{fm}$ would give
$M_1c^2\simeq1.97\ \mathrm{TeV}$, and a smaller circle gives a heavier tower. For a field this
is the whole story: the smaller the circle, the better it hides. For a string it is half of
the story, as \cref{sec:circle:winding} shows.
\end{physicsbox}

\subsection{The Kaluza--Klein gauge field and the charge of the tower}\label{sec:circle:kkgauge}

A field theory that contains gravity has more than scalars. The $D$-dimensional metric
$G_{MN}$ splits, from the $(D-1)$-dimensional point of view, into a metric
$\tilde G_{\mu\nu}$, a vector $A_\mu$ and a scalar $\sigma$, packaged as
\begin{equation}\label{eq:circle:kkansatz}
 \dd s^2=\tilde G_{\mu\nu}\,\dd x^\mu\dd x^\nu+\ee^{2\sigma}\big(\dd y+A_\mu\,\dd x^\mu\big)^2
\end{equation}
\source{Tong §8.1, eq.~(8.2), ll.~11198--11212}. All three are taken independent of $y$ when we
work at energies far below $1/R$. Each has a clear meaning.
\begin{itemize}[leftmargin=*]
\item $\ee^{\sigma}$ rescales the length of the circle: its proper circumference is
$2\pi R\,\ee^{\sigma}$. A constant shift of $\sigma$ is a change of $R$. Nothing in the
reduced action fixes the value of $\sigma$; there is no potential for it. A massless scalar of
this kind is called a \emph{modulus}\index{modulus (massless scalar)}
\source{Tong §8.1, ll.~11257--11271}. In the language of the worldsheet, $R$ labels a family of
conformal field theories connected by an exactly marginal deformation
\source{GPR §2.1, ll.~458--575}.
\item $A_\mu$ is a gauge field. A coordinate change $y\to y+\Lambda(x)$ that depends on the
position in the non-compact directions is a diffeomorphism of the $D$-dimensional space. It
leaves \eqref{eq:circle:kkansatz} of the same form, because
$\dd y+A_\mu\dd x^\mu\to\dd y+(A_\mu+\partial_\mu\Lambda)\dd x^\mu$: the shift of $y$ can be
absorbed by $A_\mu\to A_\mu-\partial_\mu\Lambda$. A position-dependent rotation of the circle
is a $U(1)$ gauge transformation\index{Kaluza--Klein gauge field}. (Tong's sign of $\Lambda$
is the opposite one; nothing depends on it.) Inserting the ansatz in the Einstein--Hilbert
action gives gravity in $D-1$ dimensions coupled to a Maxwell field
$F_{\mu\nu}=\partial_\mu A_\nu-\partial_\nu A_\mu$ and to $\sigma$ \TL\
\source{Tong §8.1, ll.~11212--11256}. We do not repeat that computation; the statement we need
is only which fields appear.
\end{itemize}
The Fourier modes of \eqref{eq:circle:fourier} are \emph{charged} under this $U(1)$. Under
$y\to y+\Lambda(x)$ the mode $\Phi_n$ picks up a phase,
\begin{equation}\label{eq:circle:kkcharge}
 \Phi_n(x)\;\to\;\ee^{\ii n\Lambda(x)/R}\,\Phi_n(x),
\end{equation}
which is the transformation law of a field of charge $n/R$. Rescaling the gauge field to
$A'_\mu=A_\mu/R$ makes the charges integers \source{Tong §8.1.1, ll.~11336--11349}. The reason
is simple: the gauge symmetry is translation along the circle, and the conserved charge of a
translation is momentum. \emph{Kaluza--Klein charge is momentum in the hidden direction.}

The closed bosonic string has two more massless fields, the two-form $B_{MN}$ and the dilaton
$\Phi$ (\cref{ch:quantum,ch:backgrounds}). The dilaton reduces to a scalar. The two-form
reduces to a two-form $B_{\mu\nu}$ and a second vector,
\begin{equation}\label{eq:circle:Atilde}
 \tilde A_\mu=B_{\mu\,25}.
\end{equation}
So the low-energy theory in 25 dimensions contains a metric, a two-form, \emph{two} $U(1)$
gauge fields $A_\mu$ and $\tilde A_\mu$, and two massless scalars $\Phi$ and $\sigma$
\source{Tong §8.1, ll.~11272--11283}. We have just seen which states carry the charge of
$A_\mu$. No mode of any field in the low-energy action is charged under $\tilde A_\mu$. The
string will supply the missing charged states.

%=====================================================================================
\section{The string on a circle: two integers}\label{sec:circle:winding}

Now let a closed string move on $\R^{1,24}\times S^1$ with
\begin{equation}\label{eq:circle:ident}
 X^{25}\sim X^{25}+2\pi R .
\end{equation}
The identification changes the string in two ways \source{Tong §8.2, ll.~11351--11372}.

\subsection{Momentum is quantized}

The centre of mass of the string behaves like a particle. Its wavefunction contains the
factor $\ee^{\ii p\,x}$, where $x$ is the centre-of-mass position on the circle and $p$ the
total momentum $p^{25}$. The points $x$ and $x+2\pi R$ are the same point, so the wavefunction
must agree there: $\ee^{2\pi\ii pR}=1$, that is
\begin{equation}\label{eq:circle:momentum}
 p=\frac nR,\qquad n\in\Z .
\end{equation}
\index{momentum quantization}This is the same integer as in \eqref{eq:circle:kkmass}. Nothing
here is special to strings.

\subsection{Winding is possible}

The second change is special to strings. A closed string is a map from the circle
$\sigma\in[0,2\pi)$ into spacetime, and so far we required
$X^M(\tau,\sigma+2\pi)=X^M(\tau,\sigma)$. With the identification \eqref{eq:circle:ident} this
is too strong for $M=25$. The string is closed if the two ends of the $\sigma$-interval are at
the same \emph{point of the circle}, and the coordinate values of one point differ by multiples
of $2\pi R$:
\begin{equation}\label{eq:circle:windingbc}
 X^{25}(\tau,\sigma+2\pi)=X^{25}(\tau,\sigma)+2\pi R\,w,\qquad w\in\Z .
\end{equation}
The integer $w$ is the \emph{winding number}\index{winding number}: the number of times the
string wraps the circle, counted with orientation (\cref{fig:circle:cylinder}). Three facts
about $w$ are worth stating before any computation.
\begin{enumerate}[leftmargin=*]
\item $w$ is topological. A closed curve on a circle cannot be unwound by a continuous
deformation; in mathematical terms $w$ labels the class of the map in
$\pi_1(S^1)=\Z$. It cannot change as the free string evolves in time, because
$X(\tau,\sigma+2\pi)-X(\tau,\sigma)$ is continuous in $\tau$ and takes values in the discrete
set $2\pi R\,\Z$.
\item In interactions, winding is conserved additively: a string with $w=1$ and a string with
$w=-1$ can join into an unwound string. It is a charge, and
\cref{sec:circle:gauge} identifies the gauge field that couples to it.
\item A wound string cannot shrink to a point. Its length is at least $2\pi R|w|$, so its
energy is at least length times tension,
\begin{equation}\label{eq:circle:windingenergy}
 M_{\text{winding}}=2\pi R|w|\cdot T=2\pi R|w|\cdot\frac{1}{2\pi\ap}=\frac{|w|R}{\ap}.
\end{equation}
This estimate is classical. The quantum mass formula below contains exactly this term.
\end{enumerate}

\begin{figure}[t]
\centering
\begin{tikzpicture}[scale=1.0,>=Stealth]
 % cylinder: axis horizontal = non-compact direction, around = compact direction
 \draw (0,1) -- (9,1);
 \draw (0,-1) -- (9,-1);
 \draw (0,0) ellipse (0.35 and 1);
 \draw (9,-1) arc (-90:90:0.35 and 1);
 \draw[dashed] (9,1) arc (90:270:0.35 and 1);
 % w = 0 string: small loop on the front face
 \draw[very thick,blue!60!black] (2.0,0.05) ellipse (0.55 and 0.4);
 \node[blue!60!black] at (2.0,-1.45) {$w=0$};
 % w = 1 string: goes once around, slightly wiggly
 \draw[very thick,red!70!black]
   (5.0,-1) .. controls (5.55,-0.6) and (5.15,-0.1) .. (5.45,0.3)
            .. controls (5.6,0.6) and (5.2,0.8) .. (5.1,1);
 \draw[very thick,red!70!black,dashed]
   (5.1,1) .. controls (4.75,0.5) and (4.85,-0.5) .. (5.0,-1);
 \node[red!70!black] at (5.1,-1.45) {$w=1$};
 % w = 2 string
 \draw[very thick,green!40!black]
   (7.2,-1) .. controls (7.7,-0.5) and (7.5,0.5) .. (7.45,1);
 \draw[very thick,green!40!black,dashed]
   (7.45,1) .. controls (7.2,0.4) and (7.35,-0.5) .. (7.6,-1);
 \draw[very thick,green!40!black]
   (7.6,-1) .. controls (8.1,-0.5) and (7.9,0.5) .. (7.8,1);
 \draw[very thick,green!40!black,dashed]
   (7.8,1) .. controls (7.45,0.4) and (7.3,-0.5) .. (7.2,-1);
 \node[green!40!black] at (7.55,-1.45) {$w=2$};
 % labels
 \draw[->] (0.2,-1.9) -- (1.8,-1.9) node[right] {$X^\mu$};
 \draw[->] (-0.75,-0.6) arc (220:140:0.9);
 \node[left] at (-0.95,0) {$X^{25}$};
 \draw[<->] (9.75,-1) -- (9.75,1);
 \node[right] at (9.8,0) {$2R$};
\end{tikzpicture}
\caption{Closed strings on $\R\times S^1$, drawn as a cylinder; only one non-compact
direction is shown. The string with $w=0$ can shrink to a point. The strings with $w=1$ and
$w=2$ cannot: they wrap the circle, and their length is at least $2\pi R|w|$. Dashed parts of
the curves lie on the far side of the cylinder.}
\label{fig:circle:cylinder}
\end{figure}

\begin{pitfallbox}[Common pitfall: ``winding is just momentum in disguise'']
Momentum $n$ and winding $w$ are independent integers with different origins. The integer
$n$ exists because the \emph{wavefunction} must be single valued; it is a quantum effect, and
a particle has it too. The integer $w$ exists because the \emph{classical configuration}
$X(\sigma)$ need only be single valued up to the identification; it is present before
quantization, and a particle does not have it. The two look alike in the final formulas, and
that is a result (\cref{ch:tduality}), not an assumption.
\end{pitfallbox}

%=====================================================================================
\section{Zero modes: \texorpdfstring{$p_L$}{pL} and \texorpdfstring{$p_R$}{pR}}
\label{sec:circle:zeromodes}

\subsection{The mode expansion with winding}

In conformal gauge every $X^M$ obeys the wave equation
$\partial_+\partial_-X^M=0$, and the general solution is a left-mover plus a right-mover,
$X=X_L(\sigma^+)+X_R(\sigma^-)$ (\cref{sec:classical:modes}). The identification
\eqref{eq:circle:ident} does not change the local physics, so the wave equation, the
oscillators and their commutators are exactly as before. Only the part of $X^{25}$ that is
linear in $\tau$ and $\sigma$ changes. Write the most general linear piece as
\[
 X^{25}(\tau,\sigma)=x+a\,\tau+b\,\sigma+\text{oscillators} .
\]
The oscillator terms $\ee^{-\ii k\sigma^\pm}$ are periodic in $\sigma$, so
$X(\tau,\sigma+2\pi)-X(\tau,\sigma)=2\pi b$, and \eqref{eq:circle:windingbc} fixes $b=wR$. For
$a$, recall that the momentum density is $\Pi=\dot X/(2\pi\ap)$ (\cref{eq:classical:momentum}), so the
total momentum is
\[
 p=\int_0^{2\pi}\dd\sigma\,\frac{\dot X^{25}}{2\pi\ap}=\frac{a}{\ap}
 \qquad\Longrightarrow\qquad a=\ap p=\frac{\ap n}{R},
\]
where the oscillators drop out of the integral and \eqref{eq:circle:momentum} was used in the
last step. Therefore
\begin{equation}\label{eq:circle:zeromode}
 X^{25}(\tau,\sigma)=x+\frac{\ap n}{R}\,\tau+wR\,\sigma+\text{oscillators}
\end{equation}
\source{Tong §8.2, ll.~11373--11380}\index{zero modes}. Now split the linear part into
functions of $\sigma^\pm=\tau\pm\sigma$. Using $\tau=\frac12(\sigma^++\sigma^-)$ and
$\sigma=\frac12(\sigma^+-\sigma^-)$,
\[
 \frac{\ap n}{R}\tau+wR\,\sigma
 =\frac12\Big(\frac{\ap n}{R}+wR\Big)\sigma^+
 +\frac12\Big(\frac{\ap n}{R}-wR\Big)\sigma^- .
\]
Comparing with the non-compact expansion \cref{eq:classical:modeexp}, where the coefficient of
$\sigma^\pm$ is $\frac12\ap p$, we see that the left- and right-movers now carry
\emph{different} momenta:
\begin{equation}\label{eq:circle:pLpR}
 \boxed{\;p_L=\frac nR+\frac{wR}{\ap},\qquad p_R=\frac nR-\frac{wR}{\ap}\;}
\end{equation}
\index{left- and right-moving momenta}and the expansion of the compact coordinate is
\begin{align}
 X_L^{25}(\sigma^+)&=\tfrac12x+\tfrac12\ap\,p_L\,\sigma^+
   +\ii\sqrt{\frac\ap2}\sum_{k\neq0}\frac1k\,\tilde\alpha_k^{25}\,\ee^{-\ii k\sigma^+},\nonumber\\
 X_R^{25}(\sigma^-)&=\tfrac12x+\tfrac12\ap\,p_R\,\sigma^-
   +\ii\sqrt{\frac\ap2}\sum_{k\neq0}\frac1k\,\alpha_k^{25}\,\ee^{-\ii k\sigma^-}
   \label{eq:circle:modeexp}
\end{align}
\TL\ \source{Tong §8.2, eq.~(8.4), ll.~11381--11420}. We write $k$ for the mode number,
because $n$ is now the momentum integer. The expansions of $X^\mu$, $\mu=0,\dots,24$, are
unchanged. Two checks of \eqref{eq:circle:pLpR}:
\begin{itemize}[leftmargin=*]
\item \emph{Dimensions.} $n/R$ is an inverse length. $wR/\ap$ is length over length$^2$, also
an inverse length.
\item \emph{Decompactification.} For $w=0$ one has $p_L=p_R=p$ and
\eqref{eq:circle:modeexp} is \cref{eq:classical:modeexp}. A string that does not wind does not
know, classically, that the direction is compact.
\end{itemize}
The inverse relations are
\begin{equation}\label{eq:circle:inverse}
 \frac nR=\frac12(p_L+p_R),\qquad \frac{wR}{\ap}=\frac12(p_L-p_R).
\end{equation}
The sum of the two chiral momenta is the ordinary momentum. The difference is the winding.
In the language of \cref{ch:classical} the zero-mode oscillators are
\begin{equation}\label{eq:circle:alpha0}
 \tilde\alpha_0^{25}=\sqrt{\frac\ap2}\,p_L,\qquad \alpha_0^{25}=\sqrt{\frac\ap2}\,p_R ,
\end{equation}
while in the non-compact directions
$\alpha_0^\mu=\tilde\alpha_0^\mu=\sqrt{\ap/2}\,p^\mu$ as before
(\cref{rem:classical:circle} anticipated this).

\begin{remark}[A first look at the dual coordinate]\label{rem:circle:dualcoord}
The splitting into $X_L$ and $X_R$ suggests a second combination,
$\tilde X\equiv X_L-X_R$\index{dual coordinate}. Its zero-mode part follows from
\eqref{eq:circle:modeexp} and \eqref{eq:circle:pLpR}:
\[
 \tilde X^{25}=\tfrac12\ap(p_L-p_R)\,\tau+\tfrac12\ap(p_L+p_R)\,\sigma+\dots
 =wR\,\tau+\frac{\ap n}{R}\,\sigma+\dots
\]
Compare with \eqref{eq:circle:zeromode}: in $\tilde X$ the winding $w$ plays the role of
momentum and $n$ the role of winding, and $\tilde X$ is periodic with period
$2\pi\ap/R$ when $X$ has period $2\pi R$. This is derived here only as algebra; its meaning is
the subject of \cref{ch:tduality,ch:doubled}.
\end{remark}

\begin{remark}[Other normalizations]\label{rem:circle:gpr}
Giveon, Porrati and Rabinovici use dimensionless momenta
\[
 p_{L,R}^{\text{GPR}}=\frac{1}{\sqrt2}\Big(\frac{\sqrt\ap}{R}\,n\pm\frac{R}{\sqrt\ap}\,w\Big)
 =\sqrt{\frac\ap2}\;p_{L,R}
\]
\source{GPR §2.2, eq.~(2.2.5), ll.~576--650}. Their $p_{L,R}$ are our zero-mode oscillators
\eqref{eq:circle:alpha0}. With this normalization $L_0$ contains $\frac12p^2$ with no factor
of $\ap$. The Lean file of \cref{sec:circle:lean} sets $\ap=1$ and uses
\eqref{eq:circle:pLpR}. When comparing formulas between sources, first check which of the
two normalizations is in use; a stray factor of $2$ in a mass formula nearly always comes
from here.
\end{remark}

\begin{physicsbox}[Physical picture: a two-component charge]
A state of the string on a circle carries a pair of integers $(n,w)$. The pair
$(p_L,p_R)$ is the same information in a basis that depends on $R$. As $R$ varies, the
integers stay fixed and the point $(p_L,p_R)$ moves in the plane. One combination does not
move at all:
\begin{equation}\label{eq:circle:norm}
 \frac\ap4\big(p_L^2-p_R^2\big)=\frac\ap4\cdot4\,\frac nR\,\frac{wR}{\ap}=nw .
\end{equation}
This indefinite quadratic form on the charges is the first appearance of the $O(1,1)$
structure that \cref{ch:narain,ch:odd} generalize to $O(d,d)$. The positive-definite
combination $\frac12(p_L^2+p_R^2)$ does depend on $R$, and it is what enters the mass.
\end{physicsbox}

%=====================================================================================
\section{The mass formula and level matching}\label{sec:circle:mass}

\subsection{Derivation}

The physical-state conditions are the same as in \cref{ch:quantum}, because they come from the
worldsheet gauge symmetry and not from the shape of spacetime:
\begin{equation}\label{eq:circle:constraints}
 (L_0-1)\ket{\text{phys}}=0,\qquad(\tilde L_0-1)\ket{\text{phys}}=0 ,
\end{equation}
with the normal-ordering constant $a=1$ of the critical dimension $D=26$, and
\begin{equation}\label{eq:circle:L0}
 L_0=\frac12\alpha_0\cdot\alpha_0+N,\qquad
 \tilde L_0=\frac12\tilde\alpha_0\cdot\tilde\alpha_0+\tilde N ,
\end{equation}
where $N$ and $\tilde N$ are the right- and left-moving levels of \cref{eq:quantum:N}. The
dot product runs over all 26 directions. The only new ingredient is that the 25th component
of $\alpha_0$ differs from that of $\tilde\alpha_0$. Using \eqref{eq:circle:alpha0},
\begin{align}
 \frac12\alpha_0\cdot\alpha_0
 &=\frac12\cdot\frac\ap2\Big(p_\mu p^\mu+p_R^2\Big)=\frac\ap4\big(-M^2+p_R^2\big),\nonumber\\
 \frac12\tilde\alpha_0\cdot\tilde\alpha_0
 &=\frac\ap4\big(-M^2+p_L^2\big).\label{eq:circle:alpha0sq}
\end{align}
Here the mass is \emph{defined} as the 25-dimensional observer defines it,
\begin{equation}\label{eq:circle:massdef}
 M^2\equiv-\sum_{\mu=0}^{24}p_\mu p^\mu ,
\end{equation}
with the sum over non-compact directions only: momentum along the circle is counted as mass,
as in \cref{sec:circle:kk}. Inserting \eqref{eq:circle:alpha0sq} in
\eqref{eq:circle:constraints} gives two expressions for the same $M^2$:
\begin{equation}\label{eq:circle:twomass}
 M^2=p_R^2+\frac4\ap(N-1),\qquad M^2=p_L^2+\frac4\ap(\tilde N-1)
\end{equation}
\TL\ \source{Tong §8.2, ll.~11421--11455}. (In lightcone gauge one chooses the lightcone
directions inside $\R^{1,24}$; then $N,\tilde N$ count the 24 transverse oscillators, of
which one is $\alpha^{25}_{-k}$, and the same two equations follow.) For $p_L=p_R=0$ these
reduce to \cref{eq:quantum:massfinal}.

\emph{Subtract} the two equations in \eqref{eq:circle:twomass}:
$0=p_L^2-p_R^2+\frac4\ap(\tilde N-N)$. With \eqref{eq:circle:norm},
$p_L^2-p_R^2=4nw/\ap$, so
\begin{equation}\label{eq:circle:levelmatching}
 \boxed{\;N-\tilde N=nw\;}
\end{equation}
\index{level matching}\TL\ \source{Tong §8.2, eq.~(8.5), ll.~11456--11464}. \emph{Add} the two
equations and divide by two. From \eqref{eq:circle:pLpR},
\[
 \frac12\big(p_L^2+p_R^2\big)
 =\frac12\Big[\Big(\frac nR+\frac{wR}\ap\Big)^2+\Big(\frac nR-\frac{wR}\ap\Big)^2\Big]
 =\frac{n^2}{R^2}+\frac{w^2R^2}{\ap^2},
\]
because the cross terms $\pm2nw/\ap$ cancel. Hence
\begin{equation}\label{eq:circle:massformula}
 \boxed{\;M^2=\frac{n^2}{R^2}+\frac{w^2R^2}{\ap^2}+\frac2\ap\big(N+\tilde N-2\big)\;}
\end{equation}
\index{mass formula!on a circle}\TL\ \source{Tong §8.2, eq.~(8.6), ll.~11465--11490}. This is
the formula announced in \eqref{eq:circle:preview}, in the book's conventions. The algebra from
\eqref{eq:circle:twomass} to \eqref{eq:circle:levelmatching} and
\eqref{eq:circle:massformula} was also confirmed by a symbolic check (sympy).

\subsection{Reading the three terms}

\begin{itemize}[leftmargin=*]
\item $n^2/R^2$ is the Kaluza--Klein term \eqref{eq:circle:kkmass}. It is large for a small
circle.
\item $w^2R^2/\ap^2$ is the square of the classical energy \eqref{eq:circle:windingenergy} of
a string stretched $|w|$ times around the circle. It is large for a \emph{large} circle. The
two terms respond to $R$ in opposite ways, and their sum cannot be made small by any choice
of $R$ once both $n$ and $w$ are non-zero: by the inequality of arithmetic and geometric
means,
\begin{equation}\label{eq:circle:amgm}
 \frac{n^2}{R^2}+\frac{w^2R^2}{\ap^2}\;\ge\;\frac{2|nw|}{\ap},
 \qquad\text{with equality iff }R^2=\ap\,|n/w| .
\end{equation}
\item $\frac2\ap(N+\tilde N-2)$ is the oscillator energy with the zero-point constant $-2$,
exactly as in flat space.
\end{itemize}
The three terms are not independent, because \eqref{eq:circle:levelmatching} ties the levels
to the charges. Eliminating $\tilde N=N-nw$,
\begin{equation}\label{eq:circle:massalt}
 \ap M^2=\Big(\frac{n\sqrt\ap}{R}-\frac{wR}{\sqrt\ap}\Big)^2+4(N-1)
        =\Big(\frac{n\sqrt\ap}{R}+\frac{wR}{\sqrt\ap}\Big)^2+4(\tilde N-1),
\end{equation}
which is \eqref{eq:circle:twomass} again and is the quickest form for computing a table.

\subsection{Why level matching changes}\label{sec:circle:whylm}

In flat space level matching $N=\tilde N$ expresses the fact that no point on a closed string
is special: the operator $L_0-\tilde L_0$ generates rigid shifts
$\sigma\to\sigma+\epsilon$, and physical states must be invariant
(\cref{sec:quantum:levelmatching}). The same is true on the circle, but the generator has a new
piece. From \eqref{eq:circle:L0} and \eqref{eq:circle:alpha0sq},
\begin{equation}\label{eq:circle:rotation}
 L_0-\tilde L_0=N-\tilde N+\frac\ap4\big(p_R^2-p_L^2\big)=N-\tilde N-nw .
\end{equation}
The term $-nw$ has a direct meaning. For a wound string, shifting $\sigma$ by $\epsilon$ moves
every point of the string along the circle by $wR\,\epsilon$, because of the term
$wR\,\sigma$ in \eqref{eq:circle:zeromode}. A translation in spacetime by a distance $\delta$
acts on a state of momentum $p=n/R$ as the phase $\ee^{\ii p\delta}=\ee^{\ii nw\epsilon}$. So
part of the $\sigma$-shift is a spacetime translation, and the phase it produces must be
compensated by the oscillators. The condition $N-\tilde N=nw$ says that the total phase
vanishes.

\begin{pitfallbox}[Common pitfall: imposing $N=\tilde N$ on a circle]
On a circle $N=\tilde N$ is required only when $nw=0$. States with both momentum and winding
\emph{must} have unequal levels. For instance $(n,w)=(1,1)$ with $N=\tilde N=0$ is not a
physical state, and its ``mass'' $\ap M^2=\ap/R^2+R^2/\ap-4$ appears nowhere in the spectrum.
The lightest state with $(n,w)=(1,1)$ has $N=1$, $\tilde N=0$.
\end{pitfallbox}

%=====================================================================================
\section{The lightest states as a function of the radius}\label{sec:circle:spectrum}

It is convenient to measure the radius in string units,
\begin{equation}\label{eq:circle:r}
 r\equiv\frac{R}{\sqrt\ap},\qquad
 \ap M^2=\frac{n^2}{r^2}+w^2r^2+2\big(N+\tilde N-2\big),\qquad N-\tilde N=nw .
\end{equation}
\Cref{tab:circle:states} lists the lightest states. Each line stands for the states with the
given $(n,w)$ and their images under $(n,w)\to(-n,-w)$, which have the same mass.
\Cref{fig:circle:spectrum} plots the same functions.

\begin{table}[t]
\centering\setlength{\extrarowheight}{2.5pt}
\caption{The lightest closed-string states on a circle, with $r=R/\sqrt\ap$. The last three
columns evaluate $\ap M^2$ at $r=\frac12,1,2$. All entries follow from
\eqref{eq:circle:r}; they were recomputed with exact rational arithmetic (symbolic check).}
\label{tab:circle:states}
\smallskip
\begin{tabular}{@{}cccclccc@{}}
\toprule
$n$ & $w$ & $N$ & $\tilde N$ & $\ap M^2$ & $r=\frac12$ & $r=1$ & $r=2$\\
\midrule
$0$ & $0$ & $0$ & $0$ & $-4$ & $-4$ & $-4$ & $-4$\\
$\pm1$ & $0$ & $0$ & $0$ & $r^{-2}-4$ & $0$ & $-3$ & $-\frac{15}{4}$\\
$0$ & $\pm1$ & $0$ & $0$ & $r^{2}-4$ & $-\frac{15}{4}$ & $-3$ & $0$\\
$\pm2$ & $0$ & $0$ & $0$ & $4r^{-2}-4$ & $12$ & $0$ & $-3$\\
$0$ & $\pm2$ & $0$ & $0$ & $4r^{2}-4$ & $-3$ & $0$ & $12$\\
$0$ & $0$ & $1$ & $1$ & $0$ & $0$ & $0$ & $0$\\
$\pm1$ & $\pm1$ & $1$ & $0$ & $(r-r^{-1})^2$ & $\frac94$ & $0$ & $\frac94$\\
$\pm1$ & $\mp1$ & $0$ & $1$ & $(r-r^{-1})^2$ & $\frac94$ & $0$ & $\frac94$\\
$\pm1$ & $0$ & $1$ & $1$ & $r^{-2}$ & $4$ & $1$ & $\frac14$\\
$0$ & $\pm1$ & $1$ & $1$ & $r^{2}$ & $\frac14$ & $1$ & $4$\\
\bottomrule
\end{tabular}
\end{table}

\begin{figure}[t]
\centering
\begin{tikzpicture}[x=3.4cm,y=0.42cm,>=Stealth]
 \draw[->] (0.3,-4.8) -- (2.85,-4.8);
 \node at (1.5,-6.7) {$r=R/\sqrt\ap$};
 \draw[->] (0.3,-4.8) -- (0.3,6.8) node[above] {$\ap M^2$};
 \foreach \x in {0.5,1,1.5,2,2.5}
   \draw (\x,-4.65) -- (\x,-4.95) node[below] {\scriptsize$\x$};
 \foreach \y in {-4,-2,0,2,4,6} \draw (0.32,\y) -- (0.28,\y) node[left] {\scriptsize$\y$};
 \draw[dotted] (1,-4.8) -- (1,6.5);
 \begin{scope}
 \clip (0.3,-4.8) rectangle (2.7,6.5);
 % tachyon
 \draw[thick] (0.3,-4) -- (2.7,-4);
 % massless level
 \draw[very thick] (0.3,0) -- (2.7,0);
 % momentum modes of the tachyon
 \draw[blue!70!black,thick,domain=0.3:2.7,samples=60] plot (\x,{1/(\x*\x)-4});
 \draw[blue!70!black,thick,dashed,domain=0.3:2.7,samples=60] plot (\x,{4/(\x*\x)-4});
 % winding modes of the tachyon
 \draw[red!70!black,thick,domain=0.3:2.7,samples=60] plot (\x,{\x*\x-4});
 \draw[red!70!black,thick,dashed,domain=0.3:2.7,samples=60] plot (\x,{4*\x*\x-4});
 % (1,1) vectors
 \draw[green!40!black,very thick,domain=0.3:2.7,samples=80] plot (\x,{(\x-1/\x)*(\x-1/\x)});
 \end{scope}
 \node[right] at (2.72,-4) {\scriptsize$(0,0)$, $N=\tilde N=0$};
 \node[blue!70!black,right] at (2.72,-3.1) {\scriptsize$(\pm1,0)$};
 \node[blue!70!black,right] at (2.2,-2.45) {\scriptsize$(\pm2,0)$};
 \node[red!70!black,anchor=north west] at (2.3,1.2) {\scriptsize$(0,\pm1)$};
 \node[red!70!black,right] at (1.6,5.9) {\scriptsize$(0,\pm2)$};
 \node[green!40!black,left] at (2.62,5.6) {\scriptsize$(\pm1,\pm1)$};
 \node[right] at (2.72,0) {\scriptsize$(0,0)$, $N=\tilde N=1$};
\end{tikzpicture}
\caption{$\ap M^2$ for the states of \cref{tab:circle:states} with $N+\tilde N\le1$, labelled
by $(n,w)$. Momentum modes (blue) become light at large $r$, winding modes (red) at small $r$.
The curve $(r-1/r)^2$ (green) belongs to the vector states with $|n|=|w|=1$; it touches
zero only at $r=1$, where the curves $(\pm2,0)$ and $(0,\pm2)$ also cross zero. The picture is
symmetric under $r\to1/r$ with blue and red exchanged; \cref{ch:tduality} explains why.}
\label{fig:circle:spectrum}
\end{figure}

Several features of the table deserve comment.

\emph{The tachyon and its tower.}\index{tachyon!on a circle} The state with all quantum numbers zero is the tachyon of
\cref{ch:quantum}, $\ap M^2=-4$. Its momentum modes $(n,0)$ have $\ap M^2=n^2/r^2-4$: they are
the 26-dimensional tachyon moving along the circle, and they are tachyonic in 25 dimensions as
long as $|n|<2r$. Its winding modes $(0,w)$ have $\ap M^2=w^2r^2-4$ and are tachyonic for
$|w|<2/r$. The bosonic string has this instability at every radius; it is a property of the
bosonic vacuum and is absent in the superstring (\cref{ch:ghosts}). We keep these states in
the table because they are useful labels of the towers, not because they are physical
particles.

\emph{The massless level.} The states with $n=w=0$ and $N=\tilde N=1$ are massless for every
$R$. They are the subject of \cref{sec:circle:gauge}.

\emph{Two limits.} As $r\to\infty$ the momentum states $(n,0)$ at fixed level become
degenerate and dense: $n/R$ turns into the continuous momentum of a non-compact direction, as
it must. The winding states become infinitely heavy and leave the spectrum. As $r\to0$ the
roles are reversed: momentum states decouple, which is the field-theory expectation, but the
\emph{winding} states become dense. A continuum of states is the signature of a non-compact
dimension. So the spectrum of a string on a very small circle looks like that of a string on
a very large one. This observation, made precise, is T-duality\index{T-duality} (\cref{ch:tduality}).

\emph{The special radius.} At $r=1$ the table contains extra massless states. Setting
$M^2=0$ in \eqref{eq:circle:r} at $r=1$ and solving with level matching gives, besides the
massless level itself, exactly the following \source{Tong §8.2.3, ll.~11561--11612}:
\begin{itemize}[leftmargin=*]
\item $N=\tilde N=0$: $n^2+w^2=4$ with $nw=0$, that is $(n,w)=(\pm2,0)$ or $(0,\pm2)$. Four
scalars.
\item $N+\tilde N=1$: $n^2+w^2=2$ with $nw=N-\tilde N=\pm1$. For $nw=+1$,
$(n,w)=\pm(1,1)$ with $N=1$, $\tilde N=0$: the states
$\alpha_{-1}^\mu\ket{0;p;n,w}$. For $nw=-1$, $(n,w)=\pm(1,-1)$ with $N=0$, $\tilde N=1$:
the states $\tilde\alpha_{-1}^\mu\ket{0;p;n,w}$. Four vectors. (Tong's text prints
``$N=1$, $\tilde N=0$'' for both pairs; level matching \eqref{eq:circle:levelmatching} and the
oscillator $\tilde\alpha_{-1}$ that Tong writes for the second pair fix the assignment given
here.)
\end{itemize}
Massless charged vectors signal a non-abelian gauge symmetry. The gauge group is enhanced
from $U(1)\times U(1)$ to $SU(2)\times SU(2)$\index{enhanced gauge symmetry} \TL. We return to
this point in \cref{ch:tduality}, where $r=1$ is identified as the self-dual radius\index{self-dual radius}.

\begin{example}[Worked example: the spectrum at $R=2\sqrt\ap$]\label{ex:circle:numbers}
Take $r=2$. We list all states with $\ap M^2\le1$, using \eqref{eq:circle:massalt}.

\emph{Level $N=\tilde N=0$}, so $nw=0$. Momentum modes have
$\ap M^2=n^2/4-4$:
\[
 \ap M^2=-4,\;-\tfrac{15}{4},\;-3,\;-\tfrac74,\;0,\;\tfrac94,\;\dots
 \qquad\text{for}\qquad |n|=0,1,2,3,4,5,\dots
\]
The modes with $|n|\le3$ are tachyonic, the modes $n=\pm4$ are massless, and $|n|\ge5$ is above our cut.
Winding modes: $\ap M^2=4w^2-4$, giving $0$ for $|w|=1$ and $12$ for $|w|=2$.

\emph{Level $N=\tilde N=1$}, again $nw=0$: $\ap M^2=n^2/4$ or $4w^2$. This gives $0$ for
$n=w=0$, $\frac14$ for $n=\pm1$ and $1$ for $n=\pm2$. The winding states start at $4$.

\emph{Levels $(N,\tilde N)=(1,0)$ or $(0,1)$}, so $nw=\pm1$ and $|n|=|w|=1$:
$\ap M^2=(2-\frac12)^2=\frac94$. Too heavy. One may also try $(N,\tilde N)=(2,0)$ with
$nw=2$: the choice $(n,w)=(2,1)$ gives $\ap M^2=(1-2)^2+4=5$ and $(1,2)$ gives
$(\frac12-4)^2+4=\frac{65}{4}$.

The massless states at $r=2$ are therefore: the level $N=\tilde N=1$, $n=w=0$; two scalars
with $(n,w)=(\pm4,0)$; and two scalars with $(n,w)=(0,\pm1)$. The last four are an instance of
the accidental massless scalars that occur at $R^2=n^2\ap/4$ and $R^2=4\ap/w^2$
\source{Tong §8.2.3, ll.~11561--11586}. In physical units, if $\sqrt\ap=10^{-34}\,$m, the
first massive Kaluza--Klein state of the graviton multiplet, $\ap M^2=\frac14$, has
$M=1/(2\sqrt\ap)$, which is $\hbar c/(2\times10^{-34}\,\mathrm m)\simeq9.9\times10^{17}$~GeV.

Repeat the exercise at $r=\frac12$ and the same list of masses appears, with the roles of $n$
and $w$ exchanged: the columns $r=\frac12$ and $r=2$ of \cref{tab:circle:states} are
permutations of each other.
\end{example}

%=====================================================================================
\section{The \texorpdfstring{$U(1)\times U(1)$}{U(1) x U(1)} gauge symmetry}
\label{sec:circle:gauge}

\subsection{Massless states at a generic radius}

At a generic radius the only massless states are those with $n=w=0$ and $N=\tilde N=1$,
\[
 \alpha^{M}_{-1}\tilde\alpha^{N}_{-1}\ket{0;p},\qquad M,N\in\{\mu,25\}.
\]
In 26 non-compact dimensions these were the graviton, the $B$-field and the dilaton. A
25-dimensional observer sorts them by how many indices point along the circle
\source{Tong §8.2.1, ll.~11491--11526}:
\begin{itemize}[leftmargin=*]
\item $\alpha^\mu_{-1}\tilde\alpha^\nu_{-1}\ket{0;p}$: a symmetric traceless tensor, an
antisymmetric tensor and a trace. These are the 25-dimensional metric $G_{\mu\nu}$, two-form
$B_{\mu\nu}$ and dilaton $\Phi$.
\item $\alpha^\mu_{-1}\tilde\alpha^{25}_{-1}\ket{0;p}$ and
$\alpha^{25}_{-1}\tilde\alpha^{\mu}_{-1}\ket{0;p}$: two vectors. The symmetric combination
comes from $G_{\mu\,25}$ and is the Kaluza--Klein gauge field $A_\mu$ of
\eqref{eq:circle:kkansatz}. The antisymmetric combination comes from $B_{\mu\,25}$ and is
$\tilde A_\mu$ of \eqref{eq:circle:Atilde}.
\item $\alpha^{25}_{-1}\tilde\alpha^{25}_{-1}\ket{0;p}$: one scalar, the fluctuation of
$G_{25\,25}$, that is the radius modulus $\sigma$.
\end{itemize}
A count of states confirms that nothing is lost. In lightcone gauge the 26-dimensional massless
level has $24\times24=576$ states. In 25 dimensions the transverse index takes $23$ values plus
the circle direction, and
\[
 \underbrace{23\times23}_{G_{\mu\nu},\,B_{\mu\nu},\,\Phi}
 +\underbrace{23+23}_{A_\mu,\;\tilde A_\mu}+\underbrace{1}_{\sigma}=529+46+1=576 .
\]
The string reproduces the massless content of the Kaluza--Klein reduction of
\cref{sec:circle:kkgauge}. The existence of the massless scalar $\sigma$ is the string's way
of saying that $R$ is a free parameter\index{moduli space!of the circle}. Its vertex operator is built from
$\partial X^{25}\bar\partial X^{25}$ (\cref{ch:cft}); adding it to the worldsheet action changes
the coefficient of the kinetic term of $X^{25}$, which is a change of $R$. In the terminology
of \source{GPR §2.1, ll.~458--575} it is a truly marginal operator, and $R$ is a coordinate on
the moduli space of the conformal field theory.

\subsection{Charges: momentum and winding}

Which states are charged under the two gauge fields? Tong computes the three-point
amplitude of one photon and two tachyon modes with charges $(n,w)$ and $(-n,-w)$, and finds
it proportional to $p_L\pm p_R$, with the upper sign for $A_\mu$ and the lower sign for
$\tilde A_\mu$ \TL\ \source{Tong §8.2.2, ll.~11527--11560}. By \eqref{eq:circle:inverse},
\begin{equation}\label{eq:circle:charges}
 \text{charge under }A_\mu\ \propto\ p_L+p_R=\frac{2n}{R},\qquad
 \text{charge under }\tilde A_\mu\ \propto\ p_L-p_R=\frac{2wR}{\ap}.
\end{equation}
\index{U(1) x U(1) gauge symmetry@$U(1)\times U(1)$ gauge symmetry}The first statement is
\eqref{eq:circle:kkcharge} again. The second can be understood without any amplitude, from
the coupling of the string to a background $B$-field\index{B-field@$B$-field!winding charge} (\cref{ch:backgrounds}). In Lorentzian
conformal gauge that coupling is
\begin{equation}\label{eq:circle:Bcoupling}
 S_B=-\frac{1}{4\pi\ap}\int\dd\tau\,\dd\sigma\;\epsilon^{\alpha\beta}
     B_{MN}(X)\,\partial_\alpha X^M\partial_\beta X^N
    =-\frac{1}{2\pi\ap}\int\dd\tau\,\dd\sigma\;B_{MN}\,\dot X^MX'^N ,
\end{equation}
with $\epsilon^{\tau\sigma}=+1$; the overall sign is a convention for $B$
\source{Tong §7.2.1, eq.~(7.8), ll.~9425--9500}. Consider a string that winds the circle $w$
times and is otherwise point-like: $X^\mu=X^\mu(\tau)$ and $X'^{25}=wR$. Only the components
$B_{\mu\,25}=-B_{25\,\mu}$ contribute, and $X'^\mu=0$ removes the term with $N=\mu$, so
\begin{equation}\label{eq:circle:windingcharge}
 S_B=-\frac{1}{2\pi\ap}\int\dd\tau\;B_{\mu\,25}\,\dot X^\mu\int_0^{2\pi}\dd\sigma\,X'^{25}
    =-\frac{wR}{\ap}\int\dd\tau\;\tilde A_\mu\,\dot X^\mu .
\end{equation}
This is the worldline action of a point particle of charge $wR/\ap$ (up to the sign
convention) in the gauge field $\tilde A_\mu$, compare Tong's eq.~(7.9). The gauge
transformation of $\tilde A_\mu$ descends from the gauge symmetry
$B\to B+\dd\lambda$ of the two-form with $\lambda=\lambda_{25}(x)\,\dd y$.

\begin{physicsbox}[Physical picture: two gauge fields, two charges]
The metric gauges translations along the circle, and the conserved charge of translations is
momentum $n$. The two-form couples to the string the way a Maxwell potential couples to a
charged particle, and when the string is wrapped on the circle it \emph{is} a charged
particle in 25 dimensions, with charge $w$. A point-particle theory has $A_\mu$ but no
states charged under $\tilde A_\mu$. String theory fills both charge lattices. The integer
vector $(n,w)\in\Z^2$ is the charge vector of $U(1)\times U(1)$; in the chiral basis
$U(1)_L\times U(1)_R$ the charges are $(p_L,p_R)$.
\end{physicsbox}

%=====================================================================================
\section{What the machine has checked}\label{sec:circle:lean}

Three Lean files touch the content of this chapter. They belong to three different layers of
the library (\cref{ch:architecture}), and they illustrate three different levels of
faithfulness to the physics. None of them contains a string, a Hilbert space or a Virasoro
operator. What they contain is the algebra of the zero-mode part of
\eqref{eq:circle:massformula}.

\subsection{Stream 1: the momenta as rational functions}

The file \path{StringTheoryFormalization/UseCases/TDualityMassSpectrum.lean}, namespace
\lean{StringTheory.UseCases.TDuality}, begins with three definitions.

\begin{leanbox}[In Lean: $p_L$, $p_R$ and the momentum part of the mass]
\begin{leancode}
def leftMomentum (n w R : ℚ) : ℚ := n / R + w * R

def rightMomentum (n w R : ℚ) : ℚ := n / R - w * R

def momentumMassSq (n w R : ℚ) : ℚ :=
  (leftMomentum n w R) ^ 2 + (rightMomentum n w R) ^ 2
\end{leancode}
\textbf{Reading.} These are \eqref{eq:circle:pLpR} in units $\ap=1$, as functions
$\Q^3\to\Q$. Three things differ from the physics. (i) The radius is a rational number, not a
positive real, and may be negative. (ii) $n$ and $w$ are rationals, not integers; integrality
plays no role in the theorems below. (iii) \lean{momentumMassSq} is $p_L^2+p_R^2$, which by
the computation above \eqref{eq:circle:massformula} equals
\[
 p_L^2+p_R^2=2\Big(\frac{n^2}{R^2}+w^2R^2\Big),
\]
\emph{twice} the momentum--winding part of $\ap M^2$. The oscillator term
$\frac2\ap(N+\tilde N-2)$ and the level-matching condition are not modelled. The file's header
comment quotes the oscillator term of \eqref{eq:circle:twomass} with the coefficient $2/\ap$
where Tong's equation has $4/\ap$. Comments are not checked by Lean, and no declaration in
the file involves that term, so no theorem is affected; a reader comparing the file with this
chapter should use \eqref{eq:circle:twomass}.
\end{leanbox}

\begin{leanbox}[In Lean: how the momenta respond to $R\to1/R$, $n\leftrightarrow w$]
\begin{leancode}
theorem tduality_fixes_left_momentum (n w R : ℚ) (hR : R ≠ 0) :
    leftMomentum w n (1 / R) = leftMomentum n w R := by ...

theorem tduality_flips_right_momentum (n w R : ℚ) (hR : R ≠ 0) :
    rightMomentum w n (1 / R) = - rightMomentum n w R := by ...

theorem tduality_invariant_mass_squared (n w R : ℚ) (hR : R ≠ 0) :
    momentumMassSq w n (1 / R) = momentumMassSq n w R := by ...
\end{leancode}
\textbf{Reading.} For all rationals $n,w$ and all rational $R\neq0$: exchanging the first two
arguments and replacing $R$ by $1/R$ leaves $n/R+wR$ unchanged, changes the sign of
$n/R-wR$, and therefore leaves the sum of squares unchanged. These are identities between
rational functions. They are the algebraic core of T-duality and are discussed in
\cref{ch:tduality}; we quote them here because they are statements about the objects defined
in this chapter. They say nothing about the oscillators, about level matching (which is
invariant for the separate reason that $nw=wn$), or about interactions.
\textbf{Proof idea.} Unfold the definitions, clear denominators with \lean{field\_simp} using
$R\neq0$, and let \lean{ring} verify the polynomial identity. The third theorem rewrites with
the first two and uses $(-x)^2=x^2$.
\end{leanbox}

A companion file in the same directory, \path{NarainLattice.lean}, restates the same two
momenta under the names \lean{pL} and \lean{pR} (its namespace is that of Stream~1 with
\lean{NarainLattice} in place of \lean{TDuality}). It proves the Lean form of
\eqref{eq:circle:norm}:
\begin{leancode}
def narainNorm (n w R : ℚ) : ℚ := (pL n w R ^ 2 - pR n w R ^ 2) / 2

theorem narain_norm_R_independent (n w R : ℚ) (hR : R ≠ 0) :
    narainNorm n w R = 2 * n * w := by ...
\end{leancode}
In units $\ap=1$ this is $\frac12(p_L^2-p_R^2)=2nw$: the indefinite form does not depend on
$R$. It is the identity behind level matching, although the file does not mention levels.
\Cref{ch:narain} takes it up.

\subsection{Stream 2: the circle as the \texorpdfstring{$d=1$}{d=1} generalized metric}

The newer library \lean{DualScaleStream2} works with real matrices of arbitrary size $d$ and
with the generalized metric\index{generalized metric!circle} $\HH(G,B)$ of \cref{ch:genmetric}. For $d=1$ and $B=0$ the
metric on the circle is the $1\times1$ matrix $G=(R^2)$ in units $\ap=1$, and
\[
 \HH=\begin{pmatrix}R^2&0\\0&R^{-2}\end{pmatrix},\qquad
 Z=\begin{pmatrix}w\\n\end{pmatrix},\qquad
 Z^{\mathsf T}\HH Z=R^2w^2+\frac{n^2}{R^2}.
\]
The order $Z=(w,n)$, winding first, is the convention of this module; it follows
Hull and Zwiebach, as the file header records. The last expression is the momentum--winding
part of \eqref{eq:circle:r}. The theorem \lean{massForm\_circle} in
\path{DualScaleStream2/DFT/GeneralizedMetric.lean} states this, not for a formula written afresh,
but for the \emph{Stream~1 definition itself}.

\begin{leanbox}[In Lean: the bridge between the two libraries]
\begin{leancode}
noncomputable def massForm (H : Matrix (Charge d) (Charge d) ℝ)
    (Z : Charge d → ℝ) : ℝ :=
  Z ⬝ᵥ (H *ᵥ Z)

theorem massForm_circle (n w R : ℚ) (hR : R ≠ 0) :
    massForm (genMetric (d := 1) (!![((R : ℝ) ^ 2)]) 0)
        (Sum.elim ![(w : ℝ)] ![(n : ℝ)]) =
      ((StringTheory.UseCases.TDuality.momentumMassSq n w R : ℚ) : ℝ) / 2 := by ...
\end{leancode}
\textbf{Reading.} \lean{Charge 1} is the index type \lean{Fin 1 ⊕ Fin 1} of a doubled charge
vector; \lean{Sum.elim ![w] ![n]} is the vector whose first block is $w$ and whose second
block is $n$. For rational $n,w$ and rational $R\neq0$, all cast to real numbers, the quadratic
form $Z^{\mathsf T}\HH(R^2,0)Z$ equals one half of the rational number
\lean{momentumMassSq n w R}, cast to $\R$. The factor $\frac12$ is the factor $2$ noted in the
previous box. \textbf{Why it matters.} The two libraries were written independently, with
different number systems ($\Q$ and $\R$) and different data structures (three scalars; a
block matrix acting on a sum type). A reader might worry that they formalize two different
things under the same name. This theorem removes the worry for the circle: the kernel has
checked that they define the same function of $(n,w,R)$. It does not show that either of them
is the mass of a string; that identification is the content of
\cref{sec:circle:zeromodes,sec:circle:mass} and is Tier~L.
\textbf{Proof idea.} Compute the $1\times1$ inverse $(R^2)^{-1}$ from the adjugate formula; at
$B=0$ the generalized metric is block diagonal, so the quadratic form splits into a winding
term and a momentum term; then unfold the Stream~1 definitions, push the casts through with
\lean{push\_cast}, and finish with \lean{field\_simp} and \lean{ring}.
\end{leanbox}

\subsection{The foundation layer: an arithmetic shadow}

The third file, \path{StringTheoryFoundation/Duality/T_Duality.lean}, has no imports beyond
core Lean. Its namespace is \lean{StringTheory.Foundation.Duality.T\_Duality}.

\begin{leanbox}[In Lean: a state is a pair of integers]
\begin{leancode}
structure CircleState where
  n : Int
  w : Int
  deriving Repr, DecidableEq

def tDualState (s : CircleState) : CircleState :=
  { n := s.w, w := s.n }

theorem t_duality_involution (s : CircleState) :
    tDualState (tDualState s) = s := by ...

def selfDualMassFactor (n w : Int) : Int :=
  n * n + w * w

theorem self_dual_mass_symmetry (n w : Int) :
    selfDualMassFactor n w = selfDualMassFactor w n := by ...

theorem t_duality_spectrum_invariance (n w : Int) :
    n * n + w * w = w * w + n * n := by ...
\end{leancode}
\textbf{Reading.} A \lean{CircleState} records the charge vector $(n,w)\in\Z^2$ of
\cref{sec:circle:gauge} and nothing else: no level, no momentum $p^\mu$, no radius. The first
theorem says that swapping the two components twice is the identity. The other two say that
$n^2+w^2$ is symmetric in $n$ and $w$; the last is commutativity of integer addition. The name
\lean{selfDualMassFactor} refers to the fact that at $r=1$ the momentum--winding part of
\eqref{eq:circle:r} is $n^2+w^2$. These statements are \emph{arithmetic shadows}\index{arithmetic shadow} of the
physics: true, kernel-checked, and much weaker than their names suggest. There is no radius
in the file, so ``spectrum invariance'' here cannot refer to $R\to\ap/R$.
\textbf{Proof idea.} Definitional unfolding; the integer identities are closed by
\lean{omega}.
\end{leanbox}

\subsection{Tier table and what a faithful formalization would need}

\begin{center}
\small
\begin{tabular}{@{}>{\raggedright\arraybackslash}p{0.50\textwidth}p{0.05\textwidth}
  >{\raggedright\arraybackslash}p{0.38\textwidth}@{}}
\toprule
Statement & Tier & Where\\
\midrule
KK tower $M_n^2=n^2/R^2$; fields $G,B,\Phi,A,\tilde A,\sigma$ in 25 dimensions & \TL &
 Tong §8.1\\
Mode expansion with winding; $p_{L,R}=n/R\pm wR/\ap$ & \TL & Tong §8.2\\
$M^2=p_R^2+\frac4\ap(N-1)=p_L^2+\frac4\ap(\tilde N-1)$, hence
 \eqref{eq:circle:levelmatching}, \eqref{eq:circle:massformula} & \TL & Tong §8.2\\
Charges $p_L\pm p_R$ under $A_\mu,\tilde A_\mu$ & \TL & Tong §8.2.2\\
$n/R+wR$ fixed, $n/R-wR$ negated, sum of squares invariant under
 $(n,w,R)\to(w,n,1/R)$, over $\Q$ & \TA & \lean{tduality\_*} (three theorems)\\
$\frac12(p_L^2-p_R^2)=2nw$ for all $R\neq0$, over $\Q$, $\ap=1$ & \TA &
 \lean{narain\_norm\_R\_independent}\\
$Z^{\mathsf T}\HH(R^2,0)Z=\frac12\,$\lean{momentumMassSq} & \TA & \lean{massForm\_circle}\\
$(n,w)\mapsto(w,n)$ is an involution; $n^2+w^2$ symmetric, over $\Z$ & \TA &
 \lean{t\_duality\_involution}, \lean{self\_dual\_mass\_symmetry}\\
Table~\ref{tab:circle:states}, \cref{ex:circle:numbers}, inequality \eqref{eq:circle:amgm} &
 --- & symbolic check, not Tier~A\\
``\lean{momentumMassSq} is the mass of a string state'' & \TL & identification, never
 Tier~A\\
\bottomrule
\end{tabular}
\end{center}
The Tier~A theorems in this table depend only on the three standard axioms of Lean's
classical mathematics, namely \lean{propext}, \lean{Classical.choice} and
\lean{Quot.sound}; this was confirmed with \lean{\#print axioms} for \lean{massForm\_circle}
and \lean{tduality\_invariant\_mass\_squared}.

A formalization of \eqref{eq:circle:massformula} that deserved the name would need, at the
least: (a) a Fock space for the oscillators $\alpha_k,\tilde\alpha_k$ with the level operators
$N,\tilde N$; (b) a zero-mode sector labelled by $(n,w)\in\Z^2$ and $p^\mu$, with the
operators $L_0,\tilde L_0$ of \eqref{eq:circle:L0} acting on the tensor product; (c) the
physical-state conditions \eqref{eq:circle:constraints} as hypotheses, and
\eqref{eq:circle:levelmatching}, \eqref{eq:circle:massformula} as theorems about
eigenvalues; (d) a real radius $R>0$ and an explicit $\ap$. Step (c) is, once (a) and (b)
exist, precisely the algebra of \cref{sec:circle:mass}. Steps (a) and (b) are the real work and
are not in the library today (\cref{ch:open}).

%=====================================================================================
\begin{summarybox}
\begin{itemize}[leftmargin=*]
\item A field on $\R^{1,D-2}\times S^1$ is a tower of fields with masses $|n|/R$. The metric
component $G_{\mu\,25}$ is a $U(1)$ gauge field $A_\mu$, and the mode $n$ has charge $n$. The
radius is a modulus: a massless scalar without potential.
\item A closed string on a circle carries two integers. Momentum $n$ comes from
single-valuedness of the wavefunction; winding $w$ from the boundary condition
$X(\sigma+2\pi)=X(\sigma)+2\pi Rw$. Winding is topological, conserved, and costs energy
$|w|R/\ap$.
\item The left- and right-movers carry $p_{L,R}=n/R\pm wR/\ap$. The combination
$\frac\ap4(p_L^2-p_R^2)=nw$ is independent of $R$.
\item $M^2=n^2/R^2+w^2R^2/\ap^2+\frac2\ap(N+\tilde N-2)$ with $N-\tilde N=nw$, where $M$ is
the mass seen in 25 dimensions.
\item At generic $R$ the massless states are $G_{\mu\nu},B_{\mu\nu},\Phi$, two gauge fields
$A_\mu,\tilde A_\mu$ and the modulus $\sigma$. Momentum is the charge of $A_\mu$; winding is
the charge of $\tilde A_\mu=B_{\mu\,25}$. At $R=\sqrt\ap$ further states become massless.
\item In Lean: the rational functions $n/R\pm wR$ and their behaviour under
$(n,w,R)\to(w,n,1/R)$ \TA; the identity of the Stream~1 and Stream~2 circle mass forms,
\lean{massForm\_circle} \TA. The oscillators, level matching and the identification with a
string are not formalized.
\end{itemize}
\end{summarybox}

%=====================================================================================
\section*{Exercises}
\addcontentsline{toc}{section}{Exercises}

\begin{exercise}[$\star$ Massive field on a circle]
Repeat the reduction of \cref{sec:circle:kk} for a scalar of mass $m$ in $D$ dimensions,
adding $-\frac12m^2\Phi^2$ to \eqref{eq:circle:scalaraction}. Show that
$M_n^2=m^2+n^2/R^2$. For which $m^2<0$ and $R$ are all modes with $n\neq0$ non-tachyonic?
Compare with the momentum modes of the string tachyon in \cref{tab:circle:states}.
\end{exercise}

\begin{exercise}[$\star$ The two quadratic forms]
From \eqref{eq:circle:pLpR} verify $\frac\ap4(p_L^2-p_R^2)=nw$ and
$\frac12(p_L^2+p_R^2)=n^2/R^2+w^2R^2/\ap^2$. Express both in terms of the dimensionless
momenta of \cref{rem:circle:gpr}.
\end{exercise}

\begin{exercise}[$\star$ Accidental massless scalars]
Show that a state with $N=\tilde N=0$ is massless exactly when $r=|n|/2$ with $w=0$, or
$r=2/|w|$ with $n=0$. List the radii in the interval $\frac12\le r\le2$ at which this
happens, and the states involved.
\end{exercise}

\begin{exercise}[$\star\star$ Twisted boundary conditions]
A complex scalar may obey $\Phi(x,y+2\pi R)=\ee^{2\pi\ii\theta}\Phi(x,y)$ with a fixed
$\theta\in[0,1)$. Find the Kaluza--Klein masses. For $\theta=\frac12$ show that there is no
massless mode and that the lightest mode has $M=1/(2R)$. Explain why a \emph{real} scalar
allows only $\theta\in\{0,\frac12\}$.
\end{exercise}

\begin{exercise}[$\star\star$ States with both charges are never tachyonic]
Fix $(n,w,N,\tilde N)$ with $nw\neq0$ and regard \eqref{eq:circle:massformula} as a function
of $R$. Show that it is minimized at $R^2=\ap|n/w|$, and use level matching to show that the
minimum is $\ap M^2_{\min}=4\big(\max(N,\tilde N)-1\big)\ge0$. Which states reach
$M^2=0$, and at which radius for given $(n,w)$?
\end{exercise}

\begin{exercise}[$\star\star$ Counting tachyons]
Show that the number of tachyonic states at radius $r$ is
$\#\{n\in\Z:|n|<2r\}+\#\{w\in\Z:|w|<2/r\}-1$. Evaluate it at $r=2$ and compare with
\cref{ex:circle:numbers}. At which $r$ is the number smallest, and what is it there?
\end{exercise}

\begin{exercise}[$\star\star$ Lean: the factor of two]
State and prove in Lean, in the namespace of \lean{momentumMassSq}:
\begin{leancode}
theorem momentumMassSq_eq (n w R : ℚ) (hR : R ≠ 0) :
    momentumMassSq n w R = 2 * (n ^ 2 / R ^ 2 + w ^ 2 * R ^ 2)
\end{leancode}
Hint: \lean{unfold}, then \lean{field\_simp}, then \lean{ring}. Then prove the bound
$4nw\le{}$\lean{momentumMassSq n w R}, the Lean form of \eqref{eq:circle:amgm}. Hint: first
establish $p_L^2-p_R^2=4nw$ as a \lean{have}, then call \lean{nlinarith} with the hint
\lean{sq\_nonneg (n / R - w * R)}.
\end{exercise}

\begin{exercise}[$\star\star$ Lean: level matching as a predicate]
In the setting of \path{T_Duality.lean} define
\begin{leancode}
def levelMatched (s : CircleState) (N Nt : Int) : Prop := N - Nt = s.n * s.w
\end{leancode}
and prove that \lean{levelMatched (tDualState s) N Nt ↔ levelMatched s N Nt}. Which lemma
about integers does the proof use? State in one sentence what this theorem does and does not
say about the string spectrum.
\end{exercise}

\begin{exercise}[$\star\star\star$ The circle from the general theory]
The file \path{GeneralizedMetric.lean} proves, for every $d$, the two theorems
\lean{massForm\_covariant} and \lean{tduality\_inverts\_metric}:
\[
 Z^{\mathsf T}(g^{\mathsf T}\HH g)Z=(gZ)^{\mathsf T}\HH(gZ),\qquad
 \eta\,\HH(G,0)\,\eta=\HH(G^{-1},0).
\]
Using only these two
statements and $\eta^{\mathsf T}=\eta$, derive on paper the invariance of
$R^2w^2+n^2/R^2$ under $R\to1/R$, $n\leftrightarrow w$. Then write the corresponding Lean
statement for $d=1$ and outline how \lean{massForm\_circle} would let you transport
\lean{tduality\_invariant\_mass\_squared} from $\Q$ to the matrix setting.
\end{exercise}

%=====================================================================================
\section*{Further reading}
\addcontentsline{toc}{section}{Further reading}

\begin{itemize}[leftmargin=*]
\item Tong, \emph{String Theory}, §8.1--8.2: the text this chapter follows most closely
\source{Tong §8.1--8.2.3, ll.~11183--11625}. The coupling of the string to $B_{\mu\nu}$ used
in \cref{sec:circle:gauge} is in \source{Tong §7.2.1, ll.~9425--9500}.
\item Giveon, Porrati and Rabinovici, \emph{Target space duality in string theory}: the
circle in canonical language, with dimensionless $p_{L,R}$, the conjugate momentum and the
Hamiltonian $H=L_0+\tilde L_0$ \source{GPR §2.2, ll.~576--760}; moduli as couplings of truly
marginal operators \source{GPR §2.1, ll.~458--575}.
\item Alvarez, Alvarez-Gaum\'e and Lozano, \emph{An introduction to T-duality in string
theory}: the introduction places the circle among the dualities of statistical mechanics and
field theory, and states that the equivalence $R\to\ap/R$ of the interacting theory requires
a shift of the dilaton \source{AAL §1, ll.~95--117}. This is the starting point of
\cref{ch:tduality,ch:buscher}.
\end{itemize}

